\documentclass[12pt]{article}
\usepackage[utf8]{inputenc}
\usepackage[T1]{fontenc}
\usepackage{amsmath, amssymb, amsthm}
\usepackage{siunitx}
\usepackage{tikz}
\usepackage{pgfplots}
\usepackage{geometry}
\usepackage{titlesec}
\geometry{a4paper, margin=1in}

\pgfplotsset{compat=1.18}
\usetikzlibrary{3d, arrows.meta}

% Theorem environments
\newtheorem{definition}{Definition}
\newtheorem{proposition}{Proposition}
\newtheorem{remark}{Remark}
\newtheorem{example}{Example}
\newtheorem{theorem}{Theorem}

% Paragraph formatting
\setlength{\parindent}{1.5em}
\setlength{\parskip}{1em}

% Add extra vertical space
\titleformat{\section}{\normalfont\Large\bfseries}{\thesection}{1em}{}
\titlespacing*{\section}{0pt}{2ex plus 1ex minus .2ex}{1ex plus .2ex}

\title{Chapter 01: Double integrals }
\author{Notes from prof zeghlaoui course}

\hbadness=10000
\hfuzz=\maxdimen

\newcommand{\N}{\mathbb{N}}
\newcommand{\R}{\mathbb{R}}
\newcommand{\Z}{\mathbb{Z}}
\newcommand{\Q}{\mathbb{Q}}
\newcommand{\C}{\mathbb{C}}
\newcommand{\E}{\mathbb{E}}
\newcommand{\F}{\mathbb{F}}

% --- Derivatives ---
\newcommand{\pd}[2]{\dfrac{\partial #1}{\partial #2}}
\newcommand{\pdd}[3]{\dfrac{\partial^2 #1}{\partial #2\,\partial #3}}
\newcommand{\pddx}[2]{\dfrac{\partial^2 #1}{\partial #2^2}}
\newcommand{\dd}[2]{\dfrac{\mathrm{d} #1}{\mathrm{d} #2}}
\newcommand{\ddd}[2]{\dfrac{\mathrm{d}^2 #1}{\mathrm{d} #2^2}}
\newcommand{\diff}{\,\mathrm{d}}

%% vertical vector macro
\newcommand{\vect}[1]{\begin{pmatrix}#1\end{pmatrix}}


\begin{document}

\maketitle

\vspace{2em}
\section{Reminder:}
Recall that if $f: \left[a,b\right] \to \R$ is bounded.

for any subdivision $\sigma = \{ x_0=a <x_1< x_2<...< x_{n+1} = b\}$ of  $\left[a,b\right]$

\[S(f,\sigma)= \sum_{i = 1}^{n}M_i(x_i - x_{i-1}) \quad \text{and} \quad s(f,\sigma)= \sum_{i = 1}^{n}m_i(x_i - x_{i-1})\]
\[\text{where } m_i = \inf_{x \in [x_{i-1},x_i]} f(x), \quad M_i = \sup_{x \in [x_{i-1},x_i]} f(x)\]

called respectively, the upper and the lower Darboux sums of $f$
on $\left[a, b\right]$ with respect to the subdivision $\sigma$;

$S$: the area of the region plane delimited by:
\[
\begin{cases}
  x=a \quad y = 0 \\
  x=b \quad y = f(x) \\
\end{cases}
\]

(graph showing the lower and upper of an integral)

\[s(f,\sigma) \le S \le S(f,\sigma)\quad \text{if } \sigma_1 \subset \sigma_2,\text{ then:}\]
\[s(f,\sigma_1)  \le s(f,\sigma_2) \le S \le S(f,\sigma_2) \le S(f,\sigma_1)\]

\[
S^+ = \inf (S(f,\sigma))\quad , \quad S^- = \sup (s(f,\sigma))
\]

\[
S^- \le S \le S^+
\]

\begin{definition}
  We say that $f$ is Riemann integrable on $\left[a,b\right]$ if :
  \[S^+ = S^- = S = \int_a^b f(x)dx\]
\end{definition}

\begin{definition}
  Riemann sums

  Let $\sigma$ a subdivision of $\left[a,b\right]$ and $\zeta = \{\zeta_1,\dots,\zeta_n\} \quad / \zeta_i \in \left[x_{i-1}, x_i\right]$

  \[s(f,\sigma) \le R(f,\sigma,\zeta) = \sum_{i=1}^{n}f(\zeta_i)(x_i-x_{i-1}) \le S(f,\sigma)\]

  If $f: \left[a,b\right] \to \R$ is Riemann integrable, then:

  $\exists (\sigma_n)$ a sequence of subdivision of $\left[a,b\right]$

$\exists (\zeta_n)$ associated with 
$(\sigma_n)[a,b]\quad \zeta_n = (\zeta_n^i),\ \zeta_n^i \in [x_{i-1}^n, x_i^n]$

  $\lim\limits_{n \to \infty}R(f, \sigma_n,\zeta_n) = \int_a^bf(x)dx$

\underline{In particular:}
\[\sigma_n = \{x_i^n= a + \frac{i}{n}(b-a): i\in {0,\dots,n}\}\]
\[\lim\limits_{n \to \infty} \frac{b-a}{n} \sum_{i = 1}^{n} f(\zeta_i) = \int_a^b f(x)dx\]
\[\lim\limits_{n \to \infty} \frac{b-a}{n} \sum_{i = 1}^{n}f(a+ \frac{i}{n}(b-a)) = \int_a^bf(x)dx \quad, x_i^n-x^n_{i-1} = \frac{b-a}{n}\]
\end{definition}

\section{Double integrals:}
\subsection{Definition and first properties:}

\begin{definition}
  Darboux sums- $\quad $ Riemann sums:

  Let $f: \Delta := \left[a,b\right]\times\left[c,d\right] \to \R $  be a bounded function

  1) A subdivision $\sigma$ of $\Delta$ is a set $\sigma = \{\Delta_{ij} = \left[x_{i-1}, x_i\right]\times\left[y_{j-1}, y_j\right] \}$

  with  $(x_i)_{i = 1,\dots,n}$ is a subdivision of $\left[a,b\right]$ and $(y_j)_{j = 0,\dots,m}$ a subdivision of 
  $\left[c,d\right]$ we note : $A(\Delta_{ij}) = (x_i - x_{i-1})(y_j -y_{j-1})$

  and $\delta(\sigma) = \max_{i,j} \text{diam}(\Delta_{ij})$

  2) 
  \[\forall i,j \quad M_{ij} = \sup_{(x,y)\in\Delta_{ij}}f(x,y) \quad \text{and} \quad m_{ij} = \inf_{(x,y)\in\Delta_{ij}}f(x,y)\]

  the upper Darboux sum of $f$ with respect to $\sigma$ is :
  \[S(f,\sigma) = \sum_{i=1}^n\sum_{j=1}^m M_{ij}A(\Delta_{ij})\]

  the lower Darboux sum of $f$ with respect to $\sigma$ is :
  \[s(f,\sigma) = \sum_{i=1}^n\sum_{j=1}^m m_{ij}A(\Delta_{ij})\]

  3) if we give $\zeta = \{(\zeta_i^1,\zeta_j^2) \mid \zeta_i^1 \in \left[x_{i-1},x_i\right] \text{ and }\zeta_j^2 \in \left[y_{j-1},y_j\right]\}$

  The Riemann sum with respect to $\sigma$ and $\zeta$ is :
  \[R(f,\sigma,\zeta)= \sum_{i=1}^n\sum_{j=1}^m f(\zeta_i^1,\zeta_j^2)A(\Delta_{ij})\]
  \[R(f,\sigma,\zeta)= \sum_{i=1}^n\sum_{j=1}^m f(\zeta_i^1,\zeta_j^2)(x_i - x_{i-1})(y_j - y_{j-1})\]

\end{definition}

\subsection{Example:}

\[ f(x,y) = \alpha \quad; \forall(x,y) \in \Delta\]
\[M_{ij} = m_{ij} = f(\zeta_i^1, \zeta_j^2) = \alpha\]
\[S(f,\sigma) = s(f, \sigma) = R(f,\sigma, \zeta) = \alpha(b-a)(d-c) = \alpha A(\Delta)\]

\begin{definition}
  Riemann integrability:

  With the same notations above, we put:
  \[S^+(f)= \inf(S(f,\sigma)) \quad \text{and} \quad S^-(f)= \sup (s(f,\sigma))\]

We say that $f$ is Riemann integrable on $\Delta$ if $S^+(f) = S^-(f)$
\[\Leftrightarrow \forall \epsilon > 0; \exists \sigma \quad S(f, \sigma) - s(f, \sigma) < \epsilon\]
and we denote the common value by $\iint_{\Delta}f(x,y)dxdy$
\[\Leftrightarrow \exists(\sigma_n) \text{ of subdivisions of } \Delta / \lim S (f, \sigma_n) = \lim s(f, \sigma_n)\]

\underline{Corollary:} if  $f : \Delta = \left[a,b\right]\times\left[c,d\right]$ is Riemann integrable then:

\[\lim\limits_{n \to +\infty} \frac{(b-a)(d-c)}{n^2}\sum_{i=1}^n\sum_{j =1}^n f(a + \frac{i}{n}(b-a),c + \frac{j}{n}(d-c))\]
\[ = \iint_{\Delta}f(x,y)dxdy\]

\end{definition}

\subsection{Example:}
\[f(x,y) = x^2 + y^2 \quad,\Delta = \left[0,1\right]^2= \left[0,1\right]\times\left[0,1\right]\]


\begin{definition}
  Riemann integrability on a bounded domain:

  Let $D$ be a bounded subset of $\R^2$ and $ f: D \to \R $ a bounded function.

  We say that $f$ is Riemann integrable on $D$ if for any rectangle $\Delta \supset D$, the function :
  \[
    \tilde{f}(x,y) = 
    \begin{cases}
      f(x,y) & \text{if } (x,y) \in D \\
      0 & \text{if } (x,y) \in \Delta \setminus D
    \end{cases}
  \]

  is Riemann integrable, in this case :
  \[ \iint_D f(x,y)dxdy := \iint_{\Delta} \tilde{f}(x,y) dxdy\]

\textbf{Remark:}
Any continuous function is Riemann integrable.

\subsection{Properties:}

1)\underline{Linearity:}
If $f,g: D \to \R$ are Riemann integrable then:
$\forall \alpha,\beta \in \R \quad, \alpha f + \beta g$ is Riemann integrable and 
\[\iint_D(\alpha f + \beta g)(x,y)dxdy = \alpha\iint_Df(x,y)dxdy + \beta\iint_Dg(x,y)dxdy \]

2)

If $f \ge 0$ on $D \implies \iint_D f(x,y) dxdy \ge 0$
which is the volume of $\omega = \{(x,y,z) \in \R^3 : (x,y)\in D \text{ and } 0 \le z \le f(x,y)\}$

So; if $f \ge g \implies \iint_D\left[f(x,y) - g(x,y)\right]dxdy$

$= \text{vol}\{(x,y,z) \in \R^3: (x,y) \in D \text{ and } g(x,y) \le z \le f(x,y)\}$ 

3)\underline{Chasles's relation:}

If $D = D_1 \cup D_2$ such that $D_1 \cap D_2$ is at most a curve then:
\[\iint_Df(x,y)dxdy = \iint_{D_1}f(x,y)dxdy + \iint_{D_2}f(x,y)dxdy \]
\end{definition}

\begin{definition}
  $\text{Vol}(D) = A(D) = \iint_Ddxdy$ \quad and the average of $f$ on D is:
  $\dfrac{1}{\text{Vol}(D)}\iint_Df(x,y)dxdy = \dfrac{\iint_Df(x,y)dxdy }{\iint_D dxdy }$
\end{definition}

\textbf{\underline{Remark:}}

If  $f \equiv C \quad$, the average of $f$ is $C$

\subsection{Example:}

$D$: the triangle with vertices $(0,0),(1,0),(0,1)$
(graph) 
and $f(x,y) = xy $
\[
  \tilde{f}(x,y) = 
  \begin{cases}
    xy  & \text{if } x+y \le 1 \\
    0 & \text{if } x+y > 1
  \end{cases}
\]
$\forall x,y \in \left[0,1\right] \times \left[0,1\right]$

\[\iint_D xy dxdy = \lim\limits_{n \to +\infty} \frac{(1-0)(1-0)}{n^2} \sum_{i=1}^n\sum_{j= 1}^n\tilde{f}(\frac{i}{n},\frac{j}{n})\]

\[\lim\limits_{n \to +\infty}\sum_{i =1}^n \left[\sum_{j =1}^n \frac{i}{n}\times\frac{j}{n}\right]\]
\[U_n = \frac{1}{n^2}\left[\sum_{i=1}^n\sum_{j =1}^n \frac{i}{n}\times\frac{j}{n}\right]\]
\[U_n = \frac{1}{n^4} \sum_{i = 1}^n i\left[\sum_{j=1}^{n-i}j\right] = \frac{1}{n^4}\sum_{i =1}^n i\left(\frac{(n-i)(n-i+1)}{2}\right) \]
\[U_n =  \frac{1}{2n^4} \sum_{i=1}^n i (n^2 + i^2 - 2in + n-i) \]
\[U_n =  \frac{1}{2n^4} \sum_{i = 1}^n \left[i(n^2+n) - i^2(2n+1) +i^3\right]\]
\[U_n = \frac{1}{2n^4} \left[(n^2 + n)\frac{(n+1)n}{2} -(2n + 1)\frac{n(n+1)(2n +1)}{6} + \frac{n^2(n+1)^2}{4}\right] \]
\[U_n  = \frac{1}{4}(1 + \frac{1}{n})^2 - \frac{1}{12}(1 + \frac{1}{n})(2 + \frac{1}{n}) + \frac{1}{8}(1 + \frac{1}{n})^2\]

\[\iint_D xy dxdy = \lim\limits_{n \to + \infty} U_n = \frac{1}{4} - \frac{1}{12}\cdot 2  + \frac{1}{8} = \frac{1}{24}\]

\begin{theorem}
  Fubini's theorem:

  Let $f: D \to \R$ Riemann integrable

  1) if 
  \[D  = \{(x,y) \in \R^2 : a \le x\le b \text{ and } \exists g_1,g_2: \left[a,b\right] \to \R \text{ continuous}
    \mid 
    g_1(x) \le y \le g_2(x), \forall x \in \left[a,b\right]\}
  \]

  then: $\iint_D f(x,y)dxdy = \int_a^b \left[\int_{g_1(x)}^{g_2(x)}f(x,y)dy\right]dx$

  2) If
  \[
    D = \{ (x,y) \in \R^2 : c \le y \le d \text{ and } \exists h_1,h_2: \left[c,d\right] \to \R \text{ continuous}
    \mid h_1(y) \le x \le h_2(y), \forall y \in \left[c,d\right] \}
  \]

  then: $\iint_D f(x,y)dxdy = \int_c^d \left[\int_{h_1(y)}^{h_2(y)}f(x,y)dx\right]dy$
\end{theorem}

\subsection{Remark:}
We can use the Chasles's formula to apply this theorem.

\subsection{Example 01:}
\textbf{Problem:} Show that a disk of radius $R$ centered at the origin has area $\pi R^2$.
 
\textbf{Setup Using Cartesian Coordinates (Quick Version):}
 
The disk is:
\[
D = \{(x,y) \in \mathbb{R}^2 : x^2 + y^2 \leq R^2\}
\]
 
By symmetry (four equal quadrants), we compute:
\[
\text{Area} = 4 \iint_{D_1} dxdy
\]
 
where $D_1$ is the first quadrant portion:
\[
D_1 = \{(x,y) : 0 \leq x \leq R, \, 0 \leq y \leq \sqrt{R^2-x^2}\}
\]
 
\textbf{Compute in Cartesian Coordinates:}
 
\begin{align*}
\text{Area} &= 4\int_0^R \int_0^{\sqrt{R^2-x^2}} dy \, dx \\[0.5em]
&= 4\int_0^R \sqrt{R^2-x^2} \, dx
\end{align*}
 
Using the antiderivative $\int \sqrt{R^2-x^2} \, dx = \frac{x\sqrt{R^2-x^2}}{2} + \frac{R^2}{2}\arcsin\left(\frac{x}{R}\right)$:
 
\begin{align*}
&= 4\left[\frac{x\sqrt{R^2-x^2}}{2} + \frac{R^2}{2}\arcsin\left(\frac{x}{R}\right)\right]_0^R \\[0.5em]
&= 4\left(0 + \frac{R^2}{2} \cdot \frac{\pi}{2} - 0 - 0\right) \\[0.5em]
&= 4 \cdot \frac{\pi R^2}{4} \\[0.5em]
&= \pi R^2
\end{align*}
 
\textbf{Simpler Approach: Using Polar Coordinates:}
 
In polar coordinates $(r, \theta)$ where $x = r\cos\theta$ and $y = r\sin\theta$, the disk becomes:
\[
D = \{(r,\theta) : 0 \leq r \leq R, \, 0 \leq \theta \leq 2\pi\}
\]
 
The area element transforms: $dxdy = r \, dr \, d\theta$
 
Therefore:
\begin{align*}
\text{Area} &= \iint_D dxdy \\[0.5em]
&= \int_0^{2\pi} \int_0^R r \, dr \, d\theta \\[0.5em]
&= \int_0^{2\pi} \left[\frac{r^2}{2}\right]_0^R d\theta \\[0.5em]
&= \int_0^{2\pi} \frac{R^2}{2} \, d\theta \\[0.5em]
&= \frac{R^2}{2} \cdot 2\pi \\[0.5em]
&= \pi R^2
\end{align*}
 
\boxed{\text{Area of disk} = \pi R^2}
 
\textbf{Intuition (Polar Method):} In polar coordinates, we're summing up infinitesimal rings. Each ring at radius $r$ has circumference $2\pi r$ and thickness $dr$, so its area is $2\pi r \, dr$. Integrating from $r=0$ to $r=R$ gives the total area. The factor $r$ in the Jacobian $r \, dr \, d\theta$ accounts for the fact that rings farther from the origin have larger circumference.
 
\subsection{Example 02:}
(trapezoid area proof)
\textbf{Problem:} Prove that the area of a trapezoid with parallel sides of length $A$ and $B$ and height $h$ is $\frac{(A+B)h}{2}$.
 
\textbf{Setup:} Consider a trapezoid in the $xy$-plane where:
\begin{itemize}
  \item The bottom base (at $y=0$) has length $A$
  \item The top base (at $y=h$) has length $B$
  \item The height is $h$
  \item The left side has slope $\frac{a}{h}$ (starting at $x = 0$ when $y = 0$)
  \item The right side has slope $\frac{a+B-A}{h}$ (starting at $x = A$ when $y = 0$)
\end{itemize}
 
\textbf{Domain Description:}
 
For a fixed height $y$ (where $0 \leq y \leq h$), the trapezoid spans horizontally:
\[
D = \left\{ (x,y) \in \mathbb{R}^2 \mid 0 \leq y \leq h, \quad \frac{a}{h}y \leq x \leq \frac{a+B-A}{h}y + A \right\}
\]
 
The left boundary is: $x_{\text{left}}(y) = \frac{a}{h}y$
 
The right boundary is: $x_{\text{right}}(y) = \frac{a+B-A}{h}y + A$
 
\textbf{The Width at Height $y$:}
 
The width of the trapezoid at height $y$ is the distance between left and right boundaries:
\begin{align*}
\text{Width}(y) &= x_{\text{right}}(y) - x_{\text{left}}(y) \\
&= \left(\frac{a+B-A}{h}y + A\right) - \frac{a}{h}y \\
&= \frac{B-A}{h}y + A
\end{align*}
 
\textbf{Notice:} This is linear in $y$!
\begin{itemize}
  \item At $y = 0$: Width $= A$ (the bottom base)
  \item At $y = h$: Width $= \frac{B-A}{h} \cdot h + A = B$ (the top base)
\end{itemize}
 
\textbf{Compute the Area:}
 
\begin{align*}
A(D) &= \int_0^h \left[\int_{\frac{a}{h}y}^{\frac{a+B-A}{h}y + A} dx\right] dy \\[0.5em]
&= \int_0^h \left[\frac{B-A}{h}y + A\right] dy \\[0.5em]
&= \left[Ay + \frac{B-A}{2h}y^2\right]_0^h \\[0.5em]
&= Ah + \frac{B-A}{2h} \cdot h^2 \\[0.5em]
&= Ah + \frac{(B-A)h}{2} \\[0.5em]
&= \frac{2Ah + (B-A)h}{2} \\[0.5em]
&= \frac{(A+B)h}{2}
\end{align*}
 
\boxed{\text{Area of trapezoid} = \frac{(A+B)h}{2}}
 
\textbf{Intuition:} The width of the trapezoid changes linearly from $A$ at the bottom to $B$ at the top. The average width is $\frac{A+B}{2}$, and when multiplied by the height $h$, we get the total area.
 



\subsection{Example 03:}

\textbf{Problem:} Compute $\displaystyle \iint_D xy \, dxdy$ where $D$ is the triangular region in the first quadrant bounded by $x + y = 1$.
 
\textbf{Visualizing the Domain:}
 
The region $D$ is defined by:
\[
D = \{(x,y) \in \mathbb{R}^2 : x \geq 0, \, y \geq 0, \, x + y \leq 1\}
\]
 
This is a right triangle with vertices at $(0,0)$, $(1,0)$, and $(0,1)$.
 
\textbf{Express the Domain Explicitly:}
 
To set up the integral, we choose $x$ as the outer variable:
\begin{itemize}
  \item $x$ ranges from $0$ to $1$
  \item For each fixed $x$, $y$ ranges from $0$ to $1-x$
\end{itemize}
 
Therefore:
\[
D = \{(x,y) \in \mathbb{R}^2 : 0 \leq x \leq 1 \text{ and } 0 \leq y \leq 1-x\}
\]
 
\textbf{Set Up the Double Integral:}
 
\begin{align*}
\iint_D xy \, dxdy &= \int_0^1 x \left[\int_0^{1-x} y \, dy\right] dx
\end{align*}
 
\textbf{Evaluate the Inner Integral:}
 
The inner integral integrates $y$ with respect to $y$:
\begin{align*}
\int_0^{1-x} y \, dy &= \left[\frac{y^2}{2}\right]_0^{1-x} \\
&= \frac{(1-x)^2}{2} - 0 \\
&= \frac{(1-x)^2}{2}
\end{align*}
 
So our integral becomes:
\begin{align*}
\iint_D xy \, dxdy &= \int_0^1 x \cdot \frac{(1-x)^2}{2} \, dx \\
&= \frac{1}{2}\int_0^1 x(1-x)^2 \, dx
\end{align*}
 
\textbf{Expand $(1-x)^2$:}
 
\begin{align*}
(1-x)^2 &= 1 - 2x + x^2
\end{align*}
 
Therefore:
\begin{align*}
x(1-x)^2 &= x(1 - 2x + x^2) \\
&= x - 2x^2 + x^3
\end{align*}
 
\textbf{Evaluate the Outer Integral:}
 
\begin{align*}
\frac{1}{2}\int_0^1 (x - 2x^2 + x^3) \, dx &= \frac{1}{2}\left[\frac{x^2}{2} - \frac{2x^3}{3} + \frac{x^4}{4}\right]_0^1 \\[0.5em]
&= \frac{1}{2}\left(\frac{1}{2} - \frac{2}{3} + \frac{1}{4}\right) \\[0.5em]
&= \frac{1}{2}\left(\frac{6 - 8 + 3}{12}\right) \\[0.5em]
&= \frac{1}{2} \cdot \frac{1}{12} \\[0.5em]
&= \frac{1}{24}
\end{align*}
 
\boxed{\iint_D xy \, dxdy = \frac{1}{24}}
 
\textbf{Intuition:} The function $xy$ is small near the axes (where $x$ or $y$ is near $0$) and largest near the point $(1/3, 1/3)$ inside the triangle. The integral captures this weighted sum.
 


\subsection{Example 04:}
Let us compute volume of $\Omega$ bounded above by the paraboloid $z = 1 - x^2 -y^2$ and below by the plane $z = 1-y$

Compute the volume of the solid $\Omega$ bounded above by the paraboloid $z = 1 - x^2 - y^2$ and below by the plane $z = 1 - y$.
 
 
\textbf{Step 1: Find the Projection (Domain $D$)}
 
The solid is bounded by two surfaces:
\begin{align*}
z &= 1 - x^2 - y^2 \quad \text{(paraboloid)} \\
z &= 1 - y \quad \text{(plane)}
\end{align*}
 
Setting them equal to find where they intersect:
\begin{align*}
1 - x^2 - y^2 &= 1 - y \\
x^2 + y^2 - y &= 0 \\
x^2 + \left(y - \frac{1}{2}\right)^2 &= \frac{1}{4}
\end{align*}
 
This is a circle centered at $\left(0, \frac{1}{2}\right)$ with radius $\frac{1}{2}$.
 
The domain $D$ is:
\[
D = \left\{ (x,y) \in \mathbb{R}^2 \mid x^2 + y^2 - y \leq 0 \right\}
\]
 
This describes the region inside the circle where the plane is below the paraboloid.
 
\textbf{Step 2: Set Up the Solid Region}
 
The solid $\Omega$ is described as:
\[
\Omega = \left\{ (x,y,z) \in \mathbb{R}^3 \mid (x,y) \in D \text{ and } 1-y \leq z \leq 1-x^2-y^2 \right\}
\]
 
\textbf{Step 3: Compute the Volume}
 
The volume is:
\begin{align*}
V &= \iint_D \left[(1-x^2-y^2)-(1-y)\right] \, dxdy \\
&= \iint_D (y - x^2 - y^2) \, dxdy
\end{align*}
 
\textbf{Step 4: Express the Domain Explicitly}
 
From $x^2 + y^2 - y \leq 0$, we can write $x^2 \leq y - y^2 = y(1-y)$.
 
For the region to exist, we need $y(1-y) \geq 0$, which means $0 \leq y \leq 1$.
 
For each fixed $y$, we have:
\[
-\sqrt{y - y^2} \leq x \leq \sqrt{y - y^2}
\]
 
Therefore:
\[
D = \left\{ (x,y) \in \mathbb{R}^2 \mid 0 \leq y \leq 1, \quad -\sqrt{y - y^2} \leq x \leq \sqrt{y - y^2} \right\}
\]
 
\textbf{Step 5: Evaluate the Inner Integral}
 
\begin{align*}
V &= \int_0^1 \left[\int_{-\sqrt{y-y^2}}^{\sqrt{y-y^2}} (y - y^2 - x^2) \, dx\right] dy \\
&= \int_0^1 \left[(y-y^2)x - \frac{1}{3}x^3\right]_{-\sqrt{y-y^2}}^{\sqrt{y-y^2}} dy
\end{align*}
 
By symmetry (the integrand is even in $x$):
\begin{align*}
&= \int_0^1 2\left[(y-y^2)\sqrt{y-y^2} - \frac{1}{3}\left(\sqrt{y-y^2}\right)^3\right] dy \\
&= \int_0^1 2\sqrt{y-y^2}\left[(y-y^2) - \frac{1}{3}(y-y^2)\right] dy \\
&= \int_0^1 2\sqrt{y-y^2} \cdot \frac{2}{3}(y-y^2) \, dy \\
&= \frac{4}{3}\int_0^1 (y-y^2)^{3/2} \, dy
\end{align*}
 
\textbf{Step 6: Trigonometric Substitution}
 
\textbf{Why substitute?} We have $(y - y^2)^{3/2}$. To handle this, we complete the square:
\begin{align*}
y - y^2 &= \frac{1}{4} - \left(y - \frac{1}{2}\right)^2
\end{align*}
 
This suggests a sine substitution. Let:
\begin{align*}
y - \frac{1}{2} &= \frac{1}{2}\sin(t)
\end{align*}
 
Then:
\begin{align*}
y - y^2 &= \frac{1}{4} - \frac{1}{4}\sin^2(t) = \frac{1}{4}(1 - \sin^2(t)) = \frac{1}{4}\cos^2(t) \\
\sqrt{y-y^2} &= \frac{1}{2}|\cos(t)| = \frac{1}{2}\cos(t) \quad \text{(for } t \in [-\pi/2, \pi/2])\\
dy &= \frac{1}{2}\cos(t) \, dt
\end{align*}
 
\textbf{Bounds:} 
- When $y = 0$: $\sin(t) = -1 \Rightarrow t = -\pi/2$
- When $y = 1$: $\sin(t) = 1 \Rightarrow t = \pi/2$
 
\textbf{Step 7: Transform the Integral}
 
\begin{align*}
V &= \frac{4}{3}\int_{-\pi/2}^{\pi/2} \left(\frac{1}{2}\cos(t)\right)^3 \cdot \frac{1}{2}\cos(t) \, dt \\
&= \frac{4}{3}\int_{-\pi/2}^{\pi/2} \frac{1}{16}\cos^4(t) \, dt \\
&= \frac{1}{12}\int_{-\pi/2}^{\pi/2} \cos^4(t) \, dt
\end{align*}
 
\textbf{Step 8: Evaluate $\int_{-\pi/2}^{\pi/2} \cos^4(t) \, dt$}
 
Using the power-reduction formula for even powers:
\[
\cos^4(t) = \left(\cos^2(t)\right)^2 = \left(\frac{1 + \cos(2t)}{2}\right)^2 = \frac{1 + 2\cos(2t) + \cos^2(2t)}{4}
\]
 
For $\cos^2(2t)$:
\[
\cos^2(2t) = \frac{1 + \cos(4t)}{2}
\]
 
Therefore:
\begin{align*}
\cos^4(t) &= \frac{1}{4}\left(1 + 2\cos(2t) + \frac{1 + \cos(4t)}{2}\right) \\
&= \frac{1}{4}\left(\frac{3}{2} + 2\cos(2t) + \frac{1}{2}\cos(4t)\right) \\
&= \frac{3}{8} + \frac{1}{2}\cos(2t) + \frac{1}{8}\cos(4t)
\end{align*}
 
Now integrate:
\begin{align*}
\int_{-\pi/2}^{\pi/2} \cos^4(t) \, dt &= \int_{-\pi/2}^{\pi/2} \left(\frac{3}{8} + \frac{1}{2}\cos(2t) + \frac{1}{8}\cos(4t)\right) dt \\
&= \left[\frac{3t}{8} + \frac{1}{4}\sin(2t) + \frac{1}{32}\sin(4t)\right]_{-\pi/2}^{\pi/2} \\
&= \frac{3\pi}{16} + 0 + 0 - \left(-\frac{3\pi}{16} + 0 + 0\right) \\
&= \frac{3\pi}{8}
\end{align*}
 
\textbf{Step 9: Final Answer}
 
\begin{align*}
V &= \frac{1}{12} \cdot \frac{3\pi}{8} \\
&= \frac{3\pi}{96} \\
&= \frac{\pi}{32}
\end{align*}
 
\boxed{V = \frac{\pi}{32}}

\end{document}

